% Generated draft: requires mathematical and source review.
\documentclass[11pt]{article}
\usepackage[margin=25mm]{geometry}
\usepackage{fontspec}
\setmainfont{Latin Modern Roman}
\usepackage{amsmath,amssymb,mathtools,unicode-math}
\setmathfont{Latin Modern Math}
\usepackage{enumitem,needspace,xurl,hyperref}
\setlength{\parindent}{0pt}
\setlength{\parskip}{5pt}
\setlength{\emergencystretch}{3em}
\newcommand{\N}{\mathbb{N}}
\newcommand{\Z}{\mathbb{Z}}
\newcommand{\Q}{\mathbb{Q}}
\newcommand{\R}{\mathbb{R}}
\newcommand{\C}{\mathbb{C}}
\newcommand{\F}{\mathbb{F}}
\newcommand{\A}{\mathbb{A}}
\newcommand{\PP}{\mathbb{P}}
\newcommand{\E}{\mathbb{E}}
\newcommand{\eps}{\varepsilon}
\newcommand{\dd}{\,\mathrm{d}}
\begin{document}
\section*{UnsolvedMath ID: 800805}
\textbf{Source problem code:} catalogue entry
\subsection{Definitions and mathematical statement}
Let \(M\) denote the manifold contemplated by the source, and let \(\mathscr{D}(M)\) denote its group of diffeomorphisms, with composition as the group operation. Write \(\mathcal{M}ap(M)\) for the intended semigroup of maps \(M\to M\), also under composition, containing \(\mathscr{D}(M)\). The source does not specify the dimension, geometry, or boundary conditions of \(M\), or the regularity and admissibility conditions defining \(\mathcal{M}ap(M)\); consequently this notation does not determine a particular semigroup or topology.

For a smooth curve of diffeomorphisms \(\eta_t\), its Eulerian velocity is
\[
 u_t=\partial_t\eta_t\circ\eta_t^{-1}.
\]
Arnold's geodesic framework describes fluid evolution through a suitable geodesic equation on a diffeomorphism group. The particular metric and geometric structure to be extended are not specified in this record.

\textbf{Problem (Khesin, Misiolek, Modin, 2022, Problem 4).} Is it possible to extend Arnold's geodesic framework from \(\mathscr{D}(M)\) to \(\mathcal{M}ap(M)\) so as to capture both smooth solutions of Burgers-like and compressible fluid equations and their continuations beyond the emergence of shock waves?

Such an extension is intended to recover the smooth geodesic description while accommodating post-shock evolution through maps outside the diffeomorphism group. The source leaves unspecified the equations covered, the meaning of a geodesic on the semigroup, and the notion and admissibility of a continuation after shock formation. In particular, the displayed velocity formula uses an inverse and does not itself define velocity for a noninvertible map. These unresolved choices are part of the formulation problem; the supplied record does not determine a single fully quantified existence assertion.
\subsection{Short English statement}
The problem asks whether the geometric description of fluid motion by geodesics on a diffeomorphism group can be extended to a semigroup of maps that also describes evolution after shocks form, for Burgers-like and compressible fluid equations. The permitted maps, equations, geometric structure, and post-shock solution concept remain unspecified.
\subsection{Sources}
\begin{itemize}
\item Local UnsolvedMath registry, permanent ID \(800805\), source code \texttt{catalogue entry}, title: \emph{Geometric Hydrodynamics in Open Problems --- Problem 4}. The supplied record attributes the problem to Khesin, Misiolek, and Modin, with proposed year 2022.
\item Khesin, Misiolek, Modin, \emph{Geometric Hydrodynamics in Open Problems} (2022), Problem 4. Supplied source URL: \url{https://arxiv.org/abs/2205.01143}. The supplied provenance records access on 2026-10-16 and extraction by \texttt{arxiv\_latex}; these are supplied provenance statements, not an independent source visit or verification.
\item The local record labels the problem \texttt{open}. This is a supplied registry classification, not an independently verified assessment of the literature or current mathematical status.
\end{itemize}
\end{document}
